\chapter{Ferramentas matemáticas para a Física}
\label{cap:ferramentas-matematicas}

\section{Por que começar pela matemática?}

A Física procura descrever regularidades da natureza. Para isso, transforma observações em grandezas mensuráveis e relações entre essas grandezas. A matemática não substitui o fenômeno: ela oferece uma linguagem compacta para formular um modelo, extrair consequências e comparar previsões com medidas.

Considere uma pessoa que caminha em linha reta. A frase ``ela se afasta da origem mantendo o mesmo ritmo'' pode ser representada por

\begin{equation}
  x(t)=x_0+vt,
  \label{eq:posicao-mu}
\end{equation}

em que $x$ é a posição, $t$ é o instante, $x_0$ é a posição inicial e $v$ é a velocidade constante. A equação não é o movimento real; é um modelo que conserva apenas as características relevantes para a pergunta investigada.

Neste capítulo, revisaremos as ferramentas matemáticas que aparecerão repetidamente no livro. O objetivo não é apresentar um curso completo de matemática, mas construir um vocabulário operacional para interpretar equações, gráficos e variações.

\section{Potências e raízes}

Uma potência representa uma multiplicação repetida. Para um expoente inteiro positivo $n$,

\begin{equation}
  a^n=\underbrace{a\cdot a\cdot a\cdots a}_{n\text{ fatores}}.
\end{equation}

Por exemplo,

\begin{equation}
  3^4=3\cdot3\cdot3\cdot3=81.
\end{equation}

Para uma mesma base não nula, valem as propriedades

\begin{align}
  a^m a^n &= a^{m+n},\\
  \frac{a^m}{a^n} &= a^{m-n},\\
  \left(a^m\right)^n &= a^{mn},\\
  a^{-n} &= \frac{1}{a^n},\\
  a^0 &= 1.
\end{align}

Potências fracionárias representam raízes. Para $a\geq0$,

\begin{equation}
  a^{1/2}=\sqrt{a}
\end{equation}

e


\begin{equation}
  a^{1/n}=\sqrt[n]{a}.
\end{equation}

Essas propriedades aparecem frequentemente na manipulação de fórmulas físicas. Se uma grandeza depende de $x^2$, por exemplo, multiplicar $x$ por 3 multiplica essa grandeza por $3^2=9$.

\subsection{Potências em Python}

Em Python, o operador de potência é formado por dois asteriscos, 	exttt{**}. O símbolo 	exttt{\^} não representa potência: ele executa uma operação binária chamada \emph{XOR}.

\lstinputlisting[
  caption={Potências, raízes e uma aplicação à energia cinética.},
  label={lst:potencias-escalas}
]{codigos/cap01/potencias_e_escalas.py}

O valor calculado para $64^{1/3}$ pode aparecer como $3{,}9999999999999996$ em vez de 4. Isso ocorre porque muitos números reais não possuem representação binária finita no computador. A formatação 	exttt{:.2f} controla a apresentação do resultado, mas não altera o valor armazenado.

No trecho final do \cref{lst:potencias-escalas}, triplicamos a velocidade e verificamos numericamente que a energia cinética é multiplicada por nove. O programa, portanto, não apenas calcula um valor: ele testa uma previsão obtida pela análise matemática.

\section{Notação científica e ordem de grandeza}

A Física estuda fenômenos em escalas muito diferentes. O diâmetro aproximado de um próton é da ordem de $10^{-15}\,\si{\meter}$, enquanto a distância entre estrelas pode chegar a muitos trilhões de quilômetros. Escrever números tão pequenos ou tão grandes na forma decimal dificulta cálculos e comparações.

A notação científica representa um número na forma

\begin{equation}
  N=a\times10^n,
  \qquad 1\leq |a|<10,
\end{equation}

em que $n$ é um número inteiro.

Por exemplo,

\begin{equation}
  \SI{0.0000048}{\second}
  =\SI{4.8e-6}{\second}.
\end{equation}

O expoente $-6$ informa que a vírgula foi deslocada seis posições para a direita para obter o número $4{,}8$.

Para um número grande,

\begin{equation}
  \SI{384000000}{\meter}
  =\SI{3.84e8}{\meter}.
\end{equation}

Nesse caso, a vírgula foi deslocada oito posições para a esquerda.

\subsection{Operações com notação científica}

Na multiplicação, multiplicamos os coeficientes e somamos os expoentes:

\begin{equation}
  \left(a\times10^m\right)\left(b\times10^n\right)
  =ab\times10^{m+n}.
\end{equation}

Como exemplo,

\begin{equation}
  \left(3\times10^4\right)\left(2\times10^{-2}\right)
  =6\times10^2.
\end{equation}

Na divisão, dividimos os coeficientes e subtraímos os expoentes:

\begin{equation}
  \frac{a\times10^m}{b\times10^n}
  =\frac{a}{b}\times10^{m-n}.
\end{equation}

Portanto,

\begin{equation}
  \frac{8\times10^7}{2\times10^3}
  =4\times10^4.
\end{equation}

Na adição e na subtração, os expoentes precisam ser iguais. Por exemplo,

\begin{align}
  3{,}2\times10^5+7{,}0\times10^4
  &=3{,}2\times10^5+0{,}7\times10^5\\
  &=3{,}9\times10^5.
\end{align}

\subsection{Ordem de grandeza}

A ordem de grandeza é a potência de dez mais próxima do valor considerado. Ela permite comparar escalas e avaliar rapidamente se um resultado é plausível.

Considere

\begin{equation}
  N=a\times10^n,
  \qquad 1\leq a<10.
\end{equation}

Como o ponto médio entre $10^n$ e $10^{n+1}$, em escala logarítmica, corresponde a $\sqrt{10}\approx3{,}16$, adotamos

\begin{equation}
  \operatorname{OG}(N)=
  \begin{cases}
    10^n, & a<\sqrt{10},\\
    10^{n+1}, & a\geq\sqrt{10}.
  \end{cases}
\end{equation}

Assim, $2{,}4\times10^6$ tem ordem de grandeza $10^6$, enquanto $7{,}1\times10^6$ tem ordem de grandeza $10^7$.

\subsection{Estimativas físicas}

Antes de aceitar o resultado de um cálculo, devemos perguntar:

\begin{itemize}
  \item A unidade está correta?
  \item O sinal possui significado físico?
  \item O valor está na escala esperada?
  \item O resultado respeita os limites do modelo?
\end{itemize}

Uma velocidade cotidiana de $3\times10^7\,\si{\meter\per\second}$ é matematicamente possível como número, mas incompatível com uma caminhada, um automóvel ou um avião. Um resultado assim provavelmente indica erro de unidade, entrada ou interpretação.

\subsection*{Exemplo resolvido}

Um sinal percorre $1{,}5\times10^8\,\si{\meter}$ em $5{,}0\times10^{-1}\,\si{\second}$. Sua velocidade média é

\begin{equation}
  v_m=\frac{\Delta s}{\Delta t}.
\end{equation}

Substituindo,

\begin{align}
  v_m
  &=\frac{1{,}5\times10^8}{5{,}0\times10^{-1}}\\
  &=\frac{1{,}5}{5{,}0}\times10^{8-(-1)}\\
  &=0{,}30\times10^9\\
  &=3{,}0\times10^8\,\si{\meter\per\second}.
\end{align}

Além de efetuar o cálculo, devemos interpretar o resultado. O valor está na ordem de grandeza da velocidade da luz no vácuo, indicando que o modelo pode estar descrevendo a propagação de um sinal eletromagnético.

\subsection*{Verificação conceitual}

\begin{enumerate}
  \item Escreva $0{,}00000032\,\si{\meter}$ em notação científica.
  \item Escreva $6{,}02\times10^{23}$ na forma decimal.
  \item Determine a ordem de grandeza de $2{,}8\times10^5$.
  \item Determine a ordem de grandeza de $8{,}4\times10^5$.
  \item Explique por que a ordem de grandeza é útil para validar resultados físicos.
\end{enumerate}

\section{Proporcionalidade e leis de escala}

As relações de proporcionalidade ajudam a prever como uma grandeza física muda quando outra é alterada. Muitas vezes, essa análise pode ser feita antes de qualquer substituição numérica.

\subsection{Proporcionalidade direta}

Quando escrevemos

\begin{equation}
  y\propto x,
\end{equation}

afirmamos que existe uma constante $k$ tal que

\begin{equation}
  y=kx.
\end{equation}

A razão entre as duas grandezas permanece constante:

\begin{equation}
  \frac{y}{x}=k.
\end{equation}

Se $x$ for multiplicado por um fator $f$, então $y$ será multiplicado pelo mesmo fator:

\begin{equation}
  x'=fx
  \quad\Rightarrow\quad
  y'=fy.
\end{equation}

Por exemplo, se a distância percorrida por um objeto com velocidade constante é

\begin{equation}
  d=vt,
\end{equation}

então, mantendo $v$ constante,

\begin{equation}
  d\propto t.
\end{equation}

Se o intervalo de tempo dobrar, a distância também dobrará.

\subsection{Proporcionalidade com uma potência}

Uma relação mais geral pode ser escrita como

\begin{equation}
  y\propto x^n,
\end{equation}

ou

\begin{equation}
  y=kx^n.
\end{equation}

Quando $x$ é multiplicado por um fator $f$,

\begin{equation}
  x'=fx,
\end{equation}

temos

\begin{equation}
  y'=k(fx)^n=f^nkx^n.
\end{equation}

Portanto,

\begin{equation}
  \frac{y'}{y}=f^n.
\end{equation}

Essa expressão permite prever a mudança de $y$ sem conhecer o valor da constante $k$.

\subsection*{Exemplo resolvido: energia cinética}

A energia cinética de um corpo de massa $m$ e velocidade $v$ é

\begin{equation}
  E_c=\frac{1}{2}mv^2.
\end{equation}

Mantendo a massa constante,

\begin{equation}
  E_c\propto v^2.
\end{equation}

Se a velocidade for triplicada, $v'=3v$. Assim,

\begin{equation}
  \frac{E_c'}{E_c}
  =\left(\frac{v'}{v}\right)^2
  =3^2
  =9.
\end{equation}

Portanto, triplicar a velocidade multiplica a energia cinética por nove. Esse resultado não significa apenas que a energia aumentou. Ele mostra que a dependência quadrática torna o aumento da energia mais rápido que o aumento da velocidade.

\subsection{Proporcionalidade inversa}

Quando

\begin{equation}
  y\propto\frac{1}{x},
\end{equation}

temos

\begin{equation}
  y=\frac{k}{x}.
\end{equation}

Nesse caso, $xy=k$. Se $x$ for duplicado, $y$ será reduzido à metade.

De maneira mais geral,

\begin{equation}
  y\propto\frac{1}{x^n}.
\end{equation}

Se $x$ for multiplicado por $f$, então

\begin{equation}
  \frac{y'}{y}=\frac{1}{f^n}.
\end{equation}

\subsection*{Exemplo resolvido: intensidade e distância}

Considere um modelo idealizado no qual a intensidade $I$ recebida de uma fonte puntiforme é inversamente proporcional ao quadrado da distância $r$:

\begin{equation}
  I\propto\frac{1}{r^2}.
\end{equation}

Queremos comparar a intensidade a uma distância $r$ com a intensidade a uma distância $4r$:

\begin{equation}
  \frac{I(4r)}{I(r)}
  =\frac{1/(4r)^2}{1/r^2}
  =\frac{1}{16}.
\end{equation}

Ao quadruplicar a distância, a intensidade é reduzida a um dezesseis avos do valor inicial. Esse comportamento pode ser compreendido geometricamente: a mesma quantidade emitida pela fonte distribui-se por uma área que cresce proporcionalmente a $r^2$.

\subsection{Proporcionalidade conjunta}

Uma grandeza pode depender de mais de uma variável. Por exemplo,

\begin{equation}
  y\propto\frac{a^2b}{c^3}.
\end{equation}

Nesse caso,

\begin{equation}
  y=k\frac{a^2b}{c^3}.
\end{equation}

Se $a$ dobrar, $b$ for triplicado e $c$ permanecer constante, então

\begin{equation}
  \frac{y'}{y}=2^2\cdot3=12.
\end{equation}

A grandeza $y$ será multiplicada por 12.

\subsection{Isolamento de variáveis}

As fórmulas físicas também precisam ser reorganizadas de acordo com a grandeza procurada. Considere

\begin{equation}
  E_c=\frac{1}{2}mv^2.
\end{equation}

Para isolar $v$, multiplicamos a equação por 2:

\begin{equation}
  2E_c=mv^2.
\end{equation}

Dividimos por $m$:

\begin{equation}
  v^2=\frac{2E_c}{m}.
\end{equation}

Finalmente,

\begin{equation}
  v=\sqrt{\frac{2E_c}{m}}.
\end{equation}

Como $v$ representa o módulo da velocidade, escolhemos a raiz não negativa. Em outros contextos, uma equação quadrática pode admitir soluções positivas e negativas com significados físicos diferentes.

\subsection{Limites de uma proporcionalidade}

Uma proporcionalidade física geralmente vale dentro de determinadas condições. A expressão

\begin{equation}
  F=k\Delta x
\end{equation}

descreve aproximadamente uma mola enquanto ela permanece em seu regime elástico. Um alongamento excessivo pode deformar permanentemente o material, tornando o modelo inadequado.

Portanto, identificar uma proporcionalidade exige também perguntar:

\begin{itemize}
  \item Quais grandezas permanecem constantes?
  \item Em qual intervalo a relação foi observada?
  \item Quais hipóteses foram utilizadas?
  \item Em que situação o modelo deixa de funcionar?
\end{itemize}

\subsection*{Verificação conceitual}

\begin{enumerate}
  \item Se $y\propto x^3$, por qual fator $y$ muda quando $x$ dobra?
  \item Se $y\propto1/x^2$, por qual fator $y$ muda quando $x$ é dividido por 3?
  \item Mantendo a velocidade constante, como a energia cinética varia quando a massa é multiplicada por 5?
  \item Mantendo a massa constante, por qual fator a velocidade deve mudar para que a energia cinética seja multiplicada por 4?
  \item Se $P\propto A^2/(BC)$, determine o fator de mudança de $P$ quando $A$ dobra, $B$ triplica e $C$ é reduzido à metade.
  \item Explique por que uma proporcionalidade matematicamente correta pode deixar de representar um sistema físico real.
\end{enumerate}

\section{Logaritmos e escalas logarítmicas}

O logaritmo responde à seguinte pergunta: a qual expoente devemos elevar uma base para obter determinado número? Se

\begin{equation}
  b^y=x,
\end{equation}

então

\begin{equation}
  y=\log_bx,
\end{equation}

com as condições $b>0$, $b\neq1$ e $x>0$.

Por exemplo, $10^3=1000$ implica

\begin{equation}
  \log_{10}(1000)=3.
\end{equation}

Da mesma forma, $2^{-3}=1/8$ implica

\begin{equation}
  \log_2\left(\frac{1}{8}\right)=-3.
\end{equation}

\subsection{Bases mais utilizadas}

O logaritmo decimal possui base 10:

\begin{equation}
  \log x=\log_{10}x.
\end{equation}

O logaritmo natural possui base $\e\approx2{,}71828$ e é representado por

\begin{equation}
  \ln x=\log_\e x.
\end{equation}

O logaritmo natural aparece naturalmente em processos exponenciais, equações diferenciais, oscilações amortecidas e fenômenos de crescimento ou decaimento.

\subsection{Propriedades fundamentais}

Para $a>0$, $b>0$ e qualquer número real $n$,

\begin{align}
  \log(ab) &= \log a+\log b,\\
  \log\left(\frac{a}{b}\right) &= \log a-\log b,\\
  \log(a^n) &= n\log a.
\end{align}

Também temos $\log_b1=0$ e $\log_bb=1$. Essas propriedades transformam multiplicações em adições, divisões em subtrações e potências em produtos.

\subsection{Mudança de base}

Um logaritmo pode ser convertido de uma base para outra por meio de

\begin{equation}
  \log_bx=\frac{\log_cx}{\log_cb}.
\end{equation}

Em particular,

\begin{equation}
  \log_bx=\frac{\ln x}{\ln b}.
\end{equation}

\subsection{Logaritmos e leis de potência}

Considere uma relação física da forma

\begin{equation}
  y=Cx^n,
\end{equation}

em que $C$ e $n$ são constantes. Aplicando o logaritmo aos dois lados,

\begin{equation}
  \log y=\log C+n\log x.
\end{equation}

Definindo $Y=\log y$ e $X=\log x$, obtemos

\begin{equation}
  Y=nX+\log C.
\end{equation}

Essa é a equação de uma reta. Em um gráfico de $\log y$ em função de $\log x$, a inclinação corresponde ao expoente $n$, e a interseção com o eixo vertical corresponde a $\log C$.

\subsection*{Exemplo resolvido: determinação de uma lei de escala}

Uma grandeza $y$ apresenta os valores

\begin{center}
\begin{tabular}{cc}
\toprule
$x$ & $y$\\
\midrule
2 & 12\\
4 & 48\\
\bottomrule
\end{tabular}
\end{center}

Suponha que $y=Cx^n$. Dividindo as duas equações,

\begin{equation}
  \frac{48}{12}=\frac{C4^n}{C2^n}.
\end{equation}

Logo, $4=2^n$ e, portanto, $n=2$. Usando $x=2$ e $y=12$,

\begin{equation}
  12=C(2)^2,
\end{equation}

de onde $C=3$. Assim, a relação encontrada é

\begin{equation}
  y=3x^2.
\end{equation}

\subsection{Escalas logarítmicas}

Uma escala linear representa diferenças iguais por distâncias iguais. Uma escala logarítmica representa razões iguais por distâncias iguais. Em uma escala logarítmica de base 10, os valores $1$, $10$, $100$ e $1000$ ficam igualmente espaçados porque cada valor é dez vezes maior que o anterior.

Esse tipo de escala é útil quando uma grandeza percorre muitas ordens de magnitude. Ela permite representar valores muito diferentes no mesmo gráfico sem comprimir excessivamente os menores.

\subsection*{Exemplo físico: nível de intensidade sonora}

O nível de intensidade sonora é definido por

\begin{equation}
  \beta=10\log_{10}\left(\frac{I}{I_0}\right),
\end{equation}

em que $\beta$ é medido em decibéis, $I$ é a intensidade sonora e $I_0$ é uma intensidade de referência.

Se $I=1000I_0$, então

\begin{equation}
  \beta=10\log_{10}(1000)=\SI{30}{\decibel}.
\end{equation}

Se a intensidade for multiplicada por 10, o nível aumenta em \SI{10}{\decibel}:

\begin{align}
  \beta'
  &=10\log_{10}\left(\frac{10I}{I_0}\right)\\
  &=10\left[\log_{10}(10)+\log_{10}\left(\frac{I}{I_0}\right)\right]\\
  &=\SI{10}{\decibel}+\beta.
\end{align}

\subsection{Argumentos dimensionais}

O argumento de um logaritmo deve ser adimensional. Por isso, expressões físicas logarítmicas normalmente envolvem razões entre grandezas de mesma dimensão, como

\begin{equation}
  \ln\left(\frac{x}{x_0}\right)
\end{equation}

e

\begin{equation}
  \log_{10}\left(\frac{I}{I_0}\right).
\end{equation}

Uma expressão como $\ln(\SI{3}{\meter})$ não é fisicamente adequada, pois seu argumento ainda possui dimensão de comprimento.

\subsection{Logaritmos e exponenciais como operações inversas}

As funções exponencial e logarítmica desfazem uma à outra:

\begin{equation}
  \ln(\e^x)=x
\end{equation}

e, para $x>0$,

\begin{equation}
  \e^{\ln x}=x.
\end{equation}

\subsection*{Exemplo resolvido: tempo de decaimento}

Considere

\begin{equation}
  N(t)=N_0\e^{-t/\tau}.
\end{equation}

Queremos determinar o instante em que $N(t)=N_0/4$. Substituindo e dividindo por $N_0$,

\begin{equation}
  \frac{1}{4}=\e^{-t/\tau}.
\end{equation}

Aplicando o logaritmo natural,

\begin{equation}
  \ln\left(\frac{1}{4}\right)=-\frac{t}{\tau}.
\end{equation}

Como $\ln(1/4)=-\ln4$, obtemos

\begin{equation}
  t=\tau\ln4\approx1{,}386\tau.
\end{equation}

\subsection*{Verificação conceitual}

\begin{enumerate}
  \item Calcule $\log_{10}(10^{-6})$.
  \item Simplifique $\ln(a^3/b)$.
  \item Resolva $\e^{2x}=7$.
  \item Uma relação experimental obedece a $y=Cx^n$. Quando $x$ dobra, $y$ é multiplicado por oito. Determine $n$.
  \item Por que o argumento de um logaritmo físico deve ser adimensional?
  \item Quantos decibéis são acrescentados quando a intensidade sonora é multiplicada por 100?
  \item Em $N(t)=N_0\e^{-t/\tau}$, determine o instante em que $N=N_0/10$.
\end{enumerate}

\section{Ângulos e trigonometria}

Um ângulo pode ser medido em graus ou radianos. O radiano é a unidade natural do cálculo e relaciona o comprimento de arco $s$ ao raio $r$:

\begin{equation}
  \theta=\frac{s}{r}.
\end{equation}

Uma volta completa corresponde a $2\pi$ radianos ou $360^\circ$. Portanto,

\begin{equation}
  \theta_{\mathrm{rad}}=\theta_{\mathrm{graus}}\frac{\pi}{180}.
\end{equation}

Em um triângulo retângulo,

\begin{equation}
  \sin\theta=\frac{\text{cateto oposto}}{\text{hipotenusa}},\qquad
  \cos\theta=\frac{\text{cateto adjacente}}{\text{hipotenusa}},
\end{equation}

e


\begin{equation}
  \tan\theta=\frac{\sin\theta}{\cos\theta}.
\end{equation}

A identidade fundamental

\begin{equation}
  \sin^2\theta+\cos^2\theta=1
\end{equation}

decorre do teorema de Pitágoras aplicado ao círculo unitário.

\subsection*{Exemplo: componentes de uma força}

Uma força de módulo \SI{80}{\newton} forma um ângulo de $30^\circ$ acima da horizontal. Suas componentes cartesianas são

\begin{align}
  F_x &= F\cos\theta
       = \SI{80}{\newton}\cos 30^\circ
       \approx \SI{69.3}{\newton},\\
  F_y &= F\sin\theta
       = \SI{80}{\newton}\sin 30^\circ
       = \SI{40}{\newton}.
\end{align}

Mais adiante, representaremos essa decomposição com vetores e verificaremos graficamente que $F^2=F_x^2+F_y^2$.

\section{Funções e gráficos}

Uma função associa cada valor permitido da variável independente a um valor da variável dependente. Em Física, a notação $x(t)$ indica que a posição $x$ depende do tempo $t$.

Antes de desenhar um gráfico, devemos identificar:

\begin{enumerate}
  \item as grandezas representadas nos eixos;
  \item as unidades utilizadas;
  \item o domínio fisicamente relevante;
  \item os pontos de interseção com os eixos;
  \item o comportamento crescente, decrescente ou oscilatório;
  \item os limites do modelo.
\end{enumerate}

\subsection{Relação linear}

A função

\begin{equation}
  y=mx+b
\end{equation}

representa uma reta. O coeficiente angular

\begin{equation}
  m=\frac{\Delta y}{\Delta x}
\end{equation}

mede a variação de $y$ por unidade de $x$, enquanto $b$ é o valor de $y$ quando $x=0$.

Na \cref{eq:posicao-mu}, a inclinação de um gráfico $x\times t$ é a velocidade $v$, e a interseção com o eixo vertical é a posição inicial $x_0$. As unidades da inclinação também carregam significado:

\begin{equation}
  [v]=\frac{[x]}{[t]}=\si{\meter\per\second}.
\end{equation}

\subsection{Leis de potência e relações inversas}

Funções como $y=kx^2$ aparecem em movimentos acelerados e energias; funções como $y=k/x^2$ aparecem em modelos de campos produzidos por fontes puntiformes. A forma do gráfico ajuda a distinguir essas relações, mas a interpretação deve respeitar o domínio físico. Uma expressão que diverge em $x=0$ costuma indicar que o modelo idealizado deixou de ser aplicável nessa região.

\subsection{Comportamento exponencial}

Uma função de decaimento possui a forma

\begin{equation}
  y(t)=y_0\e^{-t/\tau},
\end{equation}

onde $\tau$ define a escala temporal. Em $t=\tau$,

\begin{equation}
  y(\tau)=\frac{y_0}{\e}\approx 0{,}368y_0.
\end{equation}

Esse comportamento surge quando a taxa de variação de uma grandeza é proporcional ao próprio valor da grandeza.

\section{Derivada: taxa de variação instantânea}

A taxa média de variação de $f$ entre $x$ e $x+\Delta x$ é

\begin{equation}
  \frac{\Delta f}{\Delta x}
  =\frac{f(x+\Delta x)-f(x)}{\Delta x}.
\end{equation}

Ao considerar intervalos progressivamente menores, definimos a derivada:

\begin{equation}
  \frac{\dd f}{\dd x}
  =\lim_{\Delta x\to 0}
  \frac{f(x+\Delta x)-f(x)}{\Delta x}.
\end{equation}

Geometricamente, a derivada é a inclinação da reta tangente ao gráfico. Fisicamente, seu significado depende das grandezas envolvidas. Para a posição $x(t)$,

\begin{equation}
  v(t)=\frac{\dd x}{\dd t},
\end{equation}

e para a velocidade,

\begin{equation}
  a(t)=\frac{\dd v}{\dd t}
      =\frac{\dd^2x}{\dd t^2}.
\end{equation}

Algumas derivadas frequentes são

\begin{align}
  \frac{\dd}{\dd x}x^n &= nx^{n-1},\\
  \frac{\dd}{\dd x}\sin x &= \cos x,\\
  \frac{\dd}{\dd x}\cos x &= -\sin x,\\
  \frac{\dd}{\dd x}\e^x &= \e^x,\\
  \frac{\dd}{\dd x}\ln x &= \frac{1}{x}.
\end{align}

As fórmulas trigonométricas pressupõem que $x$ esteja em radianos.

\subsection*{Exemplo: movimento não uniforme}

Considere

\begin{equation}
  x(t)=\SI{2}{\meter}
       +\left(\SI{3}{\meter\per\second}\right)t
       +\left(\SI{0.5}{\meter\per\second\squared}\right)t^2.
\end{equation}

Derivando em relação ao tempo,

\begin{equation}
  v(t)=\SI{3}{\meter\per\second}
       +\left(\SI{1}{\meter\per\second\squared}\right)t.
\end{equation}

A aceleração é constante e vale $\SI{1}{\meter\per\second\squared}$. Observe que os coeficientes possuem unidades diferentes para que todos os termos de $x(t)$ tenham dimensão de comprimento.

\section{Integral: acumulação}

A integral definida pode ser entendida como o limite de uma soma de pequenas contribuições. Se conhecemos a velocidade em função do tempo, o deslocamento entre $t_1$ e $t_2$ é

\begin{equation}
  \Delta x=\int_{t_1}^{t_2}v(t)\,\dd t.
\end{equation}

Geometricamente, essa integral corresponde à área algébrica sob o gráfico $v\times t$. Regiões abaixo do eixo temporal contribuem negativamente, o que distingue deslocamento de distância total percorrida.

Algumas integrais elementares são

\begin{align}
  \int x^n\,\dd x &= \frac{x^{n+1}}{n+1}+C,
  \qquad n\neq -1,\\
  \int \frac{1}{x}\,\dd x &= \ln|x|+C,\\
  \int \sin x\,\dd x &= -\cos x+C,\\
  \int \cos x\,\dd x &= \sin x+C,\\
  \int \e^x\,\dd x &= \e^x+C.
\end{align}

A constante $C$ aparece porque funções que diferem apenas por uma constante possuem a mesma derivada.

\subsection*{Exemplo: deslocamento com velocidade variável}

Para $v(t)=\alpha t$, onde $\alpha$ possui unidade de aceleração,

\begin{equation}
  \Delta x
  =\int_0^T \alpha t\,\dd t
  =\frac{1}{2}\alpha T^2.
\end{equation}

O mesmo resultado corresponde à área de um triângulo de base $T$ e altura $\alpha T$ no gráfico de velocidade por tempo.

\section{Máximos, mínimos e interpretação física}

Em um ponto interior onde uma função diferenciável possui máximo ou mínimo local, sua derivada é nula:

\begin{equation}
  f'(x)=0.
\end{equation}

Essa condição identifica pontos candidatos; não basta, sozinha, para classificá-los. Se

\begin{equation}
  f''(x)>0,
\end{equation}

a concavidade é voltada para cima e o ponto é um mínimo local. Se $f''(x)<0$, temos um máximo local. Também é necessário examinar fronteiras do domínio e pontos onde a derivada não existe.

\subsection*{Exemplo: altura máxima}

Uma pequena esfera é lançada verticalmente, e sua altura é descrita pelo modelo

\begin{equation}
  h(t)=\SI{1.2}{\meter}
       +\left(\SI{8}{\meter\per\second}\right)t
       -\frac{1}{2}\left(\SI{9.8}{\meter\per\second\squared}\right)t^2.
\end{equation}

O instante de altura máxima satisfaz $\dd h/\dd t=0$:

\begin{equation}
  8-9{,}8t=0
  \quad\Rightarrow\quad
  t\approx\SI{0.816}{\second}.
\end{equation}

Como $\dd^2h/\dd t^2=-\SI{9.8}{\meter\per\second\squared}$, a concavidade é negativa e o ponto é de máximo. A interpretação física coincide com o cálculo: no ponto mais alto, a velocidade vertical é momentaneamente nula.

\section{Estratégia para resolver problemas}

Uma solução organizada costuma seguir estas etapas:

\begin{enumerate}
  \item identificar a pergunta e o sistema físico;
  \item listar dados, incógnitas e unidades;
  \item desenhar uma representação quando ela ajudar;
  \item declarar hipóteses e escolher um modelo;
  \item desenvolver a solução simbolicamente antes de substituir números;
  \item verificar dimensões, sinal e ordem de grandeza;
  \item interpretar o resultado no contexto do problema;
  \item reconhecer limitações do modelo.
\end{enumerate}

Programas de computador podem automatizar operações e ampliar nossa capacidade de exploração, mas não escolhem sozinhos o sistema, as hipóteses ou o significado físico do resultado. Essa responsabilidade permanece com quem modela.

\section{Síntese}

Neste capítulo:

\begin{itemize}
  \item usamos notação científica e potências para lidar com escalas físicas;
  \item distinguimos proporcionalidade direta, quadrática e inversa;
  \item revisamos logaritmos e trigonometria;
  \item interpretamos gráficos como representações de relações entre grandezas;
  \item associamos derivadas a taxas instantâneas;
  \item associamos integrais a acumulações e áreas algébricas;
  \item usamos derivadas para localizar extremos;
  \item estabelecemos critérios para verificar resultados físicos.
\end{itemize}

\section{Exercícios propostos}

Os exercícios abaixo foram elaborados para este livro.

\begin{enumerate}
  \item Escreva em notação científica o diâmetro de \SI{0.00000032}{\meter}. Em seguida, converta-o para nanômetros.

  \item A energia armazenada em certo sistema é proporcional ao quadrado de uma variável $q$. Por qual fator a energia muda quando $q$ é multiplicada por 3?

  \item A intensidade de um sinal idealizado é inversamente proporcional ao quadrado da distância. Compare a intensidade a \SI{2}{\meter} e a \SI{6}{\meter} da fonte.

  \item Uma força de \SI{120}{\newton} forma $40^\circ$ com o eixo horizontal. Determine suas componentes cartesianas.

  \item Converta $150^\circ$ para radianos e $5\pi/6$ radianos para graus.

  \item Para $x(t)=4+2t-3t^2$, determine as funções velocidade e aceleração. Indique as unidades que os coeficientes devem possuir se $x$ estiver em metros e $t$ em segundos.

  \item A velocidade de um objeto é $v(t)=5-2t$, em unidades do SI. Calcule o deslocamento entre $t=0$ e $t=4\,\si{\second}$. Explique o sinal do resultado usando o gráfico $v\times t$.

  \item Determine os pontos críticos de $f(x)=x^3-6x^2+9x+2$ e classifique-os quando possível.

  \item A temperatura de um corpo é modelada por $T(t)=T_{\mathrm{amb}}+\Delta T_0\e^{-t/\tau}$. Mostre que a taxa de variação de $T-T_{\mathrm{amb}}$ é proporcional ao próprio excesso de temperatura.

  \item Uma reta em um gráfico posição-tempo passa pelos pontos $(\SI{2}{\second},\SI{7}{\meter})$ e $(\SI{8}{\second},\SI{25}{\meter})$. Determine sua inclinação, interprete-a fisicamente e encontre a posição em $t=0$.
\end{enumerate}
